\section{Discussion}
The replication and the simulation examine distinct uses of uncertainty adjustment. In the benchmark, shrinkage reallocates risk away from small-eigenvalue directions and raises realized Sharpe ratios. The simulated comparison asks whether a common adjustment preserves the allocation when relative signal uncertainty is known. Its positive compression loss identifies a further role for uncertainty information within the same quadratic portfolio framework. An effective scalar correction can improve the raw optimizer while still discarding differences relevant to the allocation.

The exposure-alignment result clarifies how these differences enter utility. A large cross-sectional dispersion in uncertainty does not by itself specify which positions are affected. In the simulation, losses are larger when the mean-scalar portfolio places more variance-scaled exposure on industries whose uncertainty differs from the average. This explains why the simulation reports both profile dispersion and signal-specific alignment. The relevant economic object is the interaction of the uncertainty profile, the signal and the resulting exposure.

The empirical design replaces the known profile with a relative estimate and selects its overall strength from scalar net utility. Consequently, the empirical utility difference measures the outcome of an implementable profile under a common selection rule. The simulated loss measures the value of removing the scalar restriction when the inputs are known. These are different estimands, and the historical result need not have the sign of the simulated compression loss. Estimation changes the information entering the portfolio, while trading costs and realized risk change the outcome used for evaluation.

The calibration and historical diagnostics help interpret that transition. The estimator captures a substantial part of the cross-sectional ordering in its synthetic regression setting, while its interval coverage varies across estimated-variance groups. In historical data, the underlying forecasts have low rank correlations with next-month returns and higher squared error than the historical-mean predictor. The heterogeneous penalty therefore operates on an estimated signal with this observed predictive behavior. The accounting in Figure~\ref{fig:decomposition} then shows where the historical portfolio difference arises: greater average gross returns are accompanied by greater net-return variance and trading expenditure.

Exposure policy is also economically material. At fixed risk aversion, the heterogeneous rule changes both portfolio composition and scale. The diagnostic direction correlations and gross-normalized position distances in Appendix~\ref{app:holdings} document the composition change. The volatility-targeted exercise changes the scale and adds a gross cap, leading to much higher absolute certainty equivalents for both strategies. Its smaller utility differences show why a statement about the uncertainty profile must be tied to a specified allocation and exposure policy.
